\section{Evaluation}\label{evaluation}

This chapter presents the CRISP-DM evaluation phase for the final set of fitted models. Chapter 6 explained how the OLS baseline, Random Forest, XGBoost, and tuned variants were specified and trained. The focus here is the held-out test-set evidence used to judge model quality, select the production model, and examine the error pattern of the retained model. Broader interpretation of what these results mean for Metro Cebu valuation practice is reserved for Chapter 8.

\subsection{\texorpdfstring{\textbf{Evaluation Approach}}{Evaluation Approach}}\label{evaluation-approach}

The study retained four evaluation metrics because each one described model quality from a different angle. All reported values were computed on the 299-row held-out test set after back-transforming predictions from log total property price to total property price in Philippine Peso.

Mean Absolute Percentage Error (MAPE) measured the average percentage gap between the actual and predicted total property price. In plain terms, it showed how far the model missed relative to the size of the observed price. This metric was retained because it is easy to compare across studies and is useful when the thesis needs a single percentage error summary.

Mean Absolute Error (MAE) measured the average absolute peso-denominated error, in this case in total Philippine Peso per property. It showed how large the miss was in the same units used by the back-transformed prediction target. This metric was retained because it is more concrete than MAPE when discussing the size of a typical prediction error.

Root Mean Square Error (RMSE) measured the square root of the average squared error. Because the errors were squared before averaging, RMSE gave more weight to larger misses than MAE. This metric was retained to show whether a model was especially sensitive to a smaller number of very large errors.

The coefficient of determination ($R^2$) measured how much of the variance in the observed total property price was explained by the model. It did not describe error size directly, but it remained useful because it showed whether the model captured the broader price pattern in the test set. The study kept $R^2$ together with MAPE, MAE, and RMSE so that no single metric would dominate the evaluation decision.

\begin{center}\rule{0.5\linewidth}{0.5pt}\end{center}

\subsection{\texorpdfstring{\textbf{Final Model Comparison}}{Final Model Comparison}}\label{final-model-comparison}

Table \ref{tab:model-comparison} presents the final comparison across the five evaluated model artifacts on the current post-cleanup sample. The tree-based models still performed much better than the OLS hedonic baseline on the held-out set, but the strongest overall result now came from the baseline Random Forest model, which achieved $R^2 = 0.807$, MAPE = 59.28\%, MAE = 4.95 million PHP and RMSE = 27.45 million PHP, all expressed in total property price. The confirmation-tuned Random Forest achieved the lowest MAPE at 53.00\%, but that gain came with materially worse $R^2$ and RMSE. Both XGBoost variants also trailed the baseline Random Forest on the retained overall criteria.

The OLS baseline performed much worse than the tree models on every retained metric. Its held-out $R^2$ fell to 0.083, while MAPE rose to 201.63\%. This gap showed that the interpretable hedonic baseline remained useful as a comparator, but it was not the best predictive option for the current open-market dataset.

The tuned models also failed to displace the retained baseline. Tuning improved XGBoost relative to its baseline and reduced MAPE for Random Forest, but neither tuned variant matched the baseline Random Forest on the held-out RMSE and $R^2$ balance used for deployment selection. 
\begin{table}[htbp]
\caption{Held-Out Test-Set Comparison of Final Model Artifacts}
\label{tab:model-comparison}
\begin{center}
\small
\begin{tabular}{p{0.31\textwidth} p{0.10\textwidth} p{0.13\textwidth} p{0.16\textwidth} p{0.16\textwidth}}
\toprule
Model & $R^2$ & MAPE & MAE & RMSE \\
\midrule
OLS (log-log baseline) & 0.083 & 201.63\% & 9.82M & 59.82M \\
Random Forest (baseline) & \textbf{0.807} & 59.28\% & \textbf{4.95M} & \textbf{27.45M} \\
Random Forest (tuned) & 0.680 & \textbf{53.00\%} & 5.08M & 35.34M \\
XGBoost (baseline) & 0.492 & 60.14\% & 6.32M & 44.54M \\
XGBoost (tuned) & 0.557 & 58.93\% & 6.06M & 41.58M \\
\bottomrule
\end{tabular}
\end{center}
\end{table}

\begin{center}\rule{0.5\linewidth}{0.5pt}\end{center}

\subsection{\texorpdfstring{\textbf{Best Model Selection}}{Best Model Selection}}\label{best-model-selection}

The baseline Random Forest model was selected as the final production model because it had the strongest overall held-out performance under the study's retained multi-metric rule. It achieved the highest $R^2$, the lowest MAE, and the lowest RMSE among all evaluated candidates in Table \ref{tab:model-comparison}. Although the confirmation-tuned Random Forest produced the lowest MAPE, its weaker RMSE and $R^2$ meant it was not the better deployment choice for total market-value prediction. The Streamlit application is therefore aligned with the saved artifact Models/rf\_model.pkl.

The baseline model also performed better than its tuned variants on the retained overall criteria. This suggested that the simpler starting configuration was already close to the useful range for the current training sample of 1{,}192 rows. The confirmation-tuned Random Forest increased the number of trees and reduced max\_features to 0.8, but its held-out performance still remained below the simpler baseline once RMSE and $R^2$ were considered. The tuned XGBoost also improved over the weaker XGBoost baseline, yet it still generalized materially worse than the retained Random Forest baseline. In this case, tuning did not overturn the incumbent model.

The comparison also showed a clear difference between predictive and explanatory roles. OLS remained the interpretable hedonic benchmark used for coefficient-based reference and for the Stage 2 MCRAI weighting process. The tree models, however, were much stronger at prediction. For that reason, the evaluation phase retained baseline Random Forest as the production estimator while keeping OLS as the econometric baseline.

This model choice satisfied the predictive part of the thesis objective. The study aimed to build a defensible valuation support model for the open-market residential segment, and the evaluation results showed that the baseline Random Forest artifact was the strongest saved model for that role.

\begin{center}\rule{0.5\linewidth}{0.5pt}\end{center}

\subsection{\texorpdfstring{\textbf{SHAP Feature Attribution}}{SHAP Feature Attribution}}\label{shap-feature-attribution}

SHAP was applied to the two retained baseline tree models so that their prediction logic could be described in a more transparent way. The ranking used mean absolute SHAP values computed on the held-out test set. Table \ref{tab:shap-top10} summarizes the top ten features for each baseline model.

The two baseline models shared a similar top-feature structure. Property-type dummies occupied the top positions in both---property\_type\_Single Detached and property\_type\_Condominium ranked first and second in both models, followed by bedrooms and dist\_consolacion\_m. Size variables (area\_sqm and floor\_area\_sqm), the Mandaue City dummy, the is\_vacant\_lot flag, and several distance features filled out the remaining top-ten positions, with minor ordering differences between the two models. These common rankings showed that the retained tree models relied on a similar mix of property-type segmentation, structural attributes, and polycentric locational signals.

An important result from this section is what did \emph{not} appear in the top ten. No individual mcrai\_* feature entered the top ten list for either baseline tree model. This does not mean that the Metro Cebu Residential Accessibility Index (MCRAI) variables had no role in prediction. It means only that, in the final SHAP ranking for the held-out set, the strongest marginal contributions came first from the distance, coordinate, and property-type variables.

\begin{table}[htbp]
\caption{Top Ten SHAP Features for the Retained Baseline Tree Models}
\label{tab:shap-top10}
\begin{center}
\small
\begin{tabular}{p{0.07\textwidth} p{0.37\textwidth} p{0.37\textwidth}}
\toprule
Rank & Random Forest & XGBoost \\
\midrule
1 & property\_type\_Single Detached & property\_type\_Single Detached \\
2 & property\_type\_Condominium & property\_type\_Condominium \\
3 & bedrooms & bedrooms \\
4 & dist\_consolacion\_m & dist\_consolacion\_m \\
5 & city\_Mandaue City & floor\_area\_sqm \\
6 & is\_vacant\_lot & city\_Mandaue City \\
7 & property\_type\_Vacant Lot & is\_vacant\_lot \\
8 & area\_sqm & longitude \\
9 & floor\_area\_sqm & dist\_naga\_city\_m \\
10 & dist\_cebu\_business\_park\_m & dist\_cebu\_business\_park\_m \\
\bottomrule
\end{tabular}
\end{center}
\end{table}

Figure \ref{fig:shap-rf} presents the Random Forest SHAP beeswarm plot, while Figure \ref{fig:shap-xgb} presents the XGBoost SHAP beeswarm plot. These figures were retained as explainability evidence for the final baseline models. The substantive interpretation of why these features dominated in Metro Cebu is deferred to Chapter 8.

\begin{figure}[htbp]
\centering
\includegraphics[width=0.92\textwidth]{../../EDA/shap_rf_summary.png}
\caption{SHAP Beeswarm Summary for the Baseline Random Forest Model}
\label{fig:shap-rf}
\end{figure}

\begin{figure}[htbp]
\centering
\includegraphics[width=0.92\textwidth]{../../EDA/shap_xgb_summary.png}
\caption{SHAP Beeswarm Summary for the Baseline XGBoost Model}
\label{fig:shap-xgb}
\end{figure}

\begin{center}\rule{0.5\linewidth}{0.5pt}\end{center}

\subsection{\texorpdfstring{\textbf{Residual Analysis}}{Residual Analysis}}\label{residual-analysis}

Residual analysis was centered on the retained production Random Forest model because that was the artifact carried forward to the application layer. The first diagnostic check was the actual-versus-predicted scatter plot shown in Figure \ref{fig:rf-actual-vs-pred}. The observations followed the expected positive pattern around the 45-degree reference line, but the spread still widened for high-value cases. This was consistent with the RMSE reported in Table \ref{tab:model-comparison}. The model captured the overall direction of price variation better than the competing artifacts, but it still produced a number of large misses at the upper end of the market.

\begin{figure}[htbp]
\centering
\includegraphics[width=0.82\textwidth]{../../EDA/rf_actual_vs_predicted.png}
\caption{Log-Scale Actual versus Predicted Total Property Price for the Baseline Random Forest Test Set}
\label{fig:rf-actual-vs-pred}
\end{figure}

Figure \ref{fig:xgb-actual-vs-pred} shows the equivalent scatter for the baseline XGBoost model. The wider spread at the upper end was consistent with XGBoost's weaker RMSE of 44.54 million PHP relative to the retained Random Forest baseline.

\begin{figure}[htbp]
\centering
\includegraphics[width=0.82\textwidth]{../../EDA/xgb_actual_vs_predicted.png}
\caption{Log-Scale Actual versus Predicted Total Property Price for the Baseline XGBoost Test Set. Compared with the baseline Random Forest, XGBoost showed a wider scatter at the upper end of the price range, consistent with its weaker RMSE of 44.54 million PHP.}
\label{fig:xgb-actual-vs-pred}
\end{figure}

No artifact-backed script in the current repository generated reliable residual breakdown tables by LGU or by property type for the held-out test set. For that reason, the residual discussion is limited here to the actual-versus-predicted scatter plot, which is the saved diagnostic figure directly produced in the workflow. That plot is still informative: it shows a positive overall relationship between actual and predicted values, but with visibly wider spread among higher-value cases. In practical terms, the production Random Forest model captured the broad direction of price variation while still producing several large misses at the upper end of the market.

\begin{center}\rule{0.5\linewidth}{0.5pt}\end{center}

\subsection{\texorpdfstring{\textbf{Benchmark Comparison}}{Benchmark Comparison}}\label{benchmark-comparison}

The production Random Forest result can be placed in context using two studies already discussed in the literature review. Nyanda et al. (2024) reported that boosting achieved a MAPE of 48.0\% in the dual-market Dar es Salaam setting, while Random Forest recorded about 52.7\% in the same mixed-market setup \parencite{nyanda2024}. The Metro Cebu deployed baseline Random Forest model posted a higher MAPE at 59.28\%, which is worse than both Nyanda et al.'s Random Forest benchmark and their boosting benchmark. The confirmation-tuned Random Forest reduced MAPE to 53.00\%, closer to Nyanda et al.'s Random Forest result, but it did so with materially weaker RMSE and $R^2$ than the deployed baseline. This comparison should be read cautiously because the market structure, price scale, and variable design were different across the studies. Even so, it reinforces that the Metro Cebu listing-based target remains noisy even when the retained model is the strongest overall local artifact.

Sousa et al. (2024) are useful in a different way. Their work did not provide a direct regression benchmark for the present model table, but it supported the use of large online listing datasets to reveal price segmentation that may not be visible in sparse official records \parencite{sousa2024}. That point matters for interpreting the Metro Cebu MAPE. The thesis relied on open-market listing data rather than clean transaction records. As a result, the evaluation used a noisier target that likely carried wider asking-price variation, seller strategy effects, and submarket mixing than a more standardized transaction database.

For that reason, the production model's MAPE should be described honestly. A MAPE of 59.28\% was high in absolute terms, especially compared with cleaner Organisation for Economic Co-operation and Development-style automated valuation settings or tightly controlled institutional datasets. At the same time, the result remained defensible for a noisy Philippine listing-based dataset that covered multiple LGUs, multiple residential property types, and a broad open-market price range. The evaluation phase therefore supported a balanced conclusion: the retained Random Forest model was the best available artifact in the current workflow on the retained overall criteria, but it was still subject to substantial prediction error that should be acknowledged clearly in the thesis.
